% TOP_PROBLEM: {"id": 30007367, "title": "The fugacity-dependent cylinder exponent for critical honeycomb loops", "category_id": 19, "category": {"id": 19, "name": "probability", "display_name": "Probability", "description": "Problems involving probability theory, stochastic processes, and random structures.", "slug": "probability", "order_index": 19, "created_at": "2026-07-31T15:26:25.671Z"}, "status": "open"}
% FIRST_POSTED: 2026-10-06
% ATTEMPT_STATUS: solved
% OPENAI_CLAIM: {"schema_version":1,"source_repo":"https://github.com/openai/math","source_commit":"adc7f1241b42e322a6451854ab7e4b4c146bf78a","release_date":"2026-10-06","match":"full","resolution":"proved","independently_reviewed":false,"second_pass":"agreed","families":[{"family":"237","title":"The three-quarter exponent for honeycomb self-avoiding walk","subject":"Probability and statistical mechanics","manuscripts":[{"title":"Cylinder loop weights and planar nesting","pdf_path":"preprints/Cylinder-loop-weights-and-planar-nesting-September-26-2026/main.pdf","readme_path":"preprints/Cylinder-loop-weights-and-planar-nesting-September-26-2026/README.md","date":"2026-09-26","theorem_ref":"thm:04-cylinder-exponent (sections/00_introduction.tex:56)","statement_sources":[{"path":"preprints/Cylinder-loop-weights-and-planar-nesting-September-26-2026/build/sections/00_introduction.tex","line":56,"label":"thm:04-cylinder-exponent"}]},{"title":"Signed cylinder propagation and marked polygons on the honeycomb lattice","pdf_path":"preprints/Signed-cylinder-propagation-and-marked-polygons-on-the-honeycomb-lattice-September-26-2026/main.pdf","readme_path":"preprints/Signed-cylinder-propagation-and-marked-polygons-on-the-honeycomb-lattice-September-26-2026/README.md","date":"2026-09-26","theorem_ref":"j:cylinder (sections/introduction.tex:80)","statement_sources":[{"path":"preprints/Signed-cylinder-propagation-and-marked-polygons-on-the-honeycomb-lattice-September-26-2026/build/sections/introduction.tex","line":80,"label":"j:cylinder"}]}],"lean":{"doc_path":"lean/docs/237.md","comparators":[],"declarations":[],"covers_main_theorem":false},"reasoning_trace":null}]}
\documentclass[11pt]{article}
\usepackage[margin=25mm]{geometry}
\usepackage{fontspec}
\setmainfont{Latin Modern Roman}
\usepackage{amsmath,amssymb,mathtools}
\usepackage{unicode-math}
\setmathfont{Latin Modern Math}
\usepackage{enumitem}
\usepackage{xurl,hyperref,bookmark}
\hypersetup{colorlinks=true,urlcolor=blue}
\setcounter{secnumdepth}{0}
\setlength{\parindent}{0pt}
\setlength{\parskip}{5pt}
\setlength{\emergencystretch}{3em}
\begin{document}
\section*{The fugacity-dependent cylinder exponent for critical honeycomb loops}

\subsection{OpenAI's claim}
The linked manuscripts present the following claims by OpenAI: Critical honeycomb loop partition exponents are claimed at every positive fugacity; fugacity two gives cylinder exponent $1/6$ and planar nesting exponent $1/12$; Signed propagation is claimed to recover the fugacity-two nesting exponents and macroscopic first-exit chord mass and length laws, together with a squared-polygon-length bound. Together they give a claimed affirmative answer in the matching scope.

\subsection{Mathematical statement of the claimed result}
\paragraph{Result 1.} On the balanced honeycomb cylinder of even circumference parameter \(N\), sum over finite vertex-disjoint polygon families separating its two prescribed markers. With \(\rho=(2+\sqrt2)^{-1/2}\), set \(Z_N(\chi)=\sum_{\mathcal P}\prod_{P\in\mathcal P}\chi\rho^{|P|}\), including the empty family. For each fixed \(\chi>0\), the source claims
\[
 Z_N(\chi)=N^{\tau(\chi)+o(1)},\qquad
 \tau(\chi)=\begin{cases}
 \frac16-\frac{2}{3\pi^2}\arccos(\chi/2)^2,&0<\chi\le2,\\
 \frac16+\frac{2}{3\pi^2}\operatorname{arcosh}(\chi/2)^2,&\chi\ge2.
 \end{cases}
\]
Separation includes winding loops and contractible loops enclosing exactly one marker. Use the quotient of the embedded honeycomb plane by the translation $iN\sqrt3/2$, with $N=2m$. The two markers are $0$ and $m e^{2\pi i/3}$.

\paragraph{Result 2.} For the honeycomb cylinder polygon partition \(G_m(g)\), each separating loop has weight \(g x^{|\lambda|}\), \(x=(2+\sqrt2)^{-1/2}\), with vertex-disjoint loop families. Along even \(m\), the source asserts
\[
 G_m(2)=m^{1/6+o(1)},\qquad
 -\left.\frac{d^2}{dg^2}\log G_m(g)\right|_{g=0}=O(\log m),
\]
where the latter counts weighted ordered incompatible cycle pairs, including self-pairs. Locally uniformly for complex \(\Theta\) with \(|\operatorname{Re}\Theta|<\pi\),
\[
 |G_m(2\cos\Theta)|\le m^{\operatorname{Re}(1/6-2\Theta^2/(3\pi^2))+o(1)}.
\]
For disjoint polygons surrounding the central face of a regular side-\(R\) hexagon, fugacity two gives \(\mathcal Z(R)=R^{1/12+o(1)}\). The cylinder identifies horizontal coordinates modulo $2m$, and its markers are $(0,0)$ and $(m,0)$, with $m$ even.

\subsection{Outline of the argument}
Contour sectors are centered and compared with a coercive fluctuation spectrum. Compensated quadratic integration and an electric contour integral continue the imbalance sum, determining cylinder growth; winding pressure and separating-loop bounds transfer it to planar nesting.

Boundary-balance identities define a two-puncture spin observable. Angular transfer and bosonization control the horizontal gas, including possible Jordan blocks, while a Pfaffian normalization fixes the physical amplitude; a separate tilted two-bond estimate gives polygon moment control.

\subsection{Scope relative to this problem}
The cylinder-loop-weights manuscript asserts the displayed piecewise exponent for every positive loop fugacity, including the transition at two. It therefore covers the complete question. The signed-cylinder companion treats fugacity two and additional propagation bounds; it would supply only that special case if considered alone. The all-fugacity theorem is the basis for the full match.

\subsection{Formal verification}
For family 237, the selected Lean statements give critical-strip mass and displacement exponents, bridge summability, and existence of pulled free energy. The $4/3$ bridge-length and $3/4$ spatial laws and the other polygon/chord exponents are outside these statements. No Comparator is assigned to the selected manuscript here; the family page selects other papers or supporting statements. This site has not executed the Lean code. The scope descriptions do not certify the remaining manuscript arguments.

\subsection{Status caveat}
OpenAI attributes these manuscripts to work by an internal model. Its release cautions that arguments without formalization can contain mistakes. This catalogue has not independently checked the proofs; the labels report the mathematical claims attributed to OpenAI.

\subsection{Sources}
\begin{itemize}
\item OpenAI, \emph{Cylinder loop weights and planar nesting}, 2026-09-26; citation key \texttt{OAI:Cylinder-loop-weights-and-planar-nesting-September-26-2026}. \href{https://github.com/openai/math/blob/adc7f1241b42e322a6451854ab7e4b4c146bf78a/preprints/Cylinder-loop-weights-and-planar-nesting-September-26-2026/main.pdf}{PDF (GitHub)}; \href{https://github.com/openai/math/raw/adc7f1241b42e322a6451854ab7e4b4c146bf78a/preprints/Cylinder-loop-weights-and-planar-nesting-September-26-2026/main.pdf}{PDF (direct)}; \href{https://github.com/openai/math/blob/adc7f1241b42e322a6451854ab7e4b4c146bf78a/preprints/Cylinder-loop-weights-and-planar-nesting-September-26-2026/README.md}{Manuscript page and citation}; \href{https://github.com/openai/math/blob/main/preprints/Cylinder-loop-weights-and-planar-nesting-September-26-2026/main.pdf}{Latest version}; Statement source: \href{https://github.com/openai/math/blob/adc7f1241b42e322a6451854ab7e4b4c146bf78a/preprints/Cylinder-loop-weights-and-planar-nesting-September-26-2026/build/sections/00\_introduction.tex\#L56}{thm:04-cylinder-exponent}.
\item OpenAI, \emph{Signed cylinder propagation and marked polygons on the honeycomb lattice}, 2026-09-26; citation key \texttt{OAI:Signed-cylinder-propagation-and-marked-polygons-on-the-honeycomb-lattice-September-26-2026}. \href{https://github.com/openai/math/blob/adc7f1241b42e322a6451854ab7e4b4c146bf78a/preprints/Signed-cylinder-propagation-and-marked-polygons-on-the-honeycomb-lattice-September-26-2026/main.pdf}{PDF (GitHub)}; \href{https://github.com/openai/math/raw/adc7f1241b42e322a6451854ab7e4b4c146bf78a/preprints/Signed-cylinder-propagation-and-marked-polygons-on-the-honeycomb-lattice-September-26-2026/main.pdf}{PDF (direct)}; \href{https://github.com/openai/math/blob/adc7f1241b42e322a6451854ab7e4b4c146bf78a/preprints/Signed-cylinder-propagation-and-marked-polygons-on-the-honeycomb-lattice-September-26-2026/README.md}{Manuscript page and citation}; \href{https://github.com/openai/math/blob/main/preprints/Signed-cylinder-propagation-and-marked-polygons-on-the-honeycomb-lattice-September-26-2026/main.pdf}{Latest version}; Statement source: \href{https://github.com/openai/math/blob/adc7f1241b42e322a6451854ab7e4b4c146bf78a/preprints/Signed-cylinder-propagation-and-marked-polygons-on-the-honeycomb-lattice-September-26-2026/build/sections/introduction.tex\#L80}{j:cylinder}.
\item Family 237: \href{https://github.com/openai/math/blob/adc7f1241b42e322a6451854ab7e4b4c146bf78a/lean/docs/237.md}{Lean scope}.
\item \href{https://github.com/openai/math/blob/adc7f1241b42e322a6451854ab7e4b4c146bf78a/CONTENTS.md}{OpenAI catalogue entry}.
\end{itemize}
\end{document}
